\subsect{Dataset Understanding Node (dataset\_node.py)}
The Dataset Understanding node performs the fan-out preparation by inspecting the provided CSV datasets, extracting basic schema information, and exposing a namespace for downstream analytics.

\begin{itemize}
  \item \textbf{Python mapping}: dataset\_node.py
  \item \textbf{Key schemas}: DatasetThoughts (from schemas.py), PLOTLY\_NAMESPACE\_KEY (from schemas.py), GraphState
  \item \textbf{Role}: Profile available datasets, build initial dataset schemas (shape and dtypes), and construct a sandbox namespace for analytics code execution.
  \item \textbf{How it operates}:
    \subitem Iterates over state["csv\_paths"], loads each CSV from the datasets/ directory into a pandas DataFrame, and records its shape and column types.
    \subitem Constructs a per-dataset schema dictionary and a namespace containing dataFrames and a dedicated Plotly namespace for capturing plots.
    \subitem Calls the LLM with a system prompt that instructs returning ONLY a JSON object with {
    (note: two fields) reasoning and schemas } where:
    \subitem reasoning: rationale for column selection and dataset approach
    \subitem schemas: the dataset schemas to guide downstream processing
    \subitem Validates the LLM response into a DatasetThoughts model.
  \item \textbf{Inputs}: state["question"], state["csv\_paths"]
  \item \textbf{Outputs}: {"dataset\_thoughts": DatasetThoughts, "namespace": {PLOTLY\_NAMESPACE\_KEY: {}, "dataFrames": dataFrames}}
  \item \textbf{Prompts design / structure}:
    \subitem System prompt explicitly requests a JSON with two fields: reasoning and schemas.
    \subitem User prompt includes the Question and Unfiltered Dataset Specifications serialized as JSON.
    \subitem The LLM return is parsed into the DatasetThoughts model so downstream components can rely on a typed structure.
\end{itemize}
